\chapter{Heat Stroke}
\index{Heat Stroke}

\section*{Definition and Overview}
Heat stroke is a life-threatening medical emergency in which the body's core temperature rises above 40 degrees Celsius and the normal mechanisms of temperature control fail. It is the most severe form of heat-related illness and is characterised by a high body temperature together with altered mental state, such as confusion, seizures, or unconsciousness. Without rapid cooling, heat stroke damages the brain, kidneys, heart, and other organs and can be fatal. Immediate first aid and emergency medical care dramatically improve the chances of survival.

\section*{Epidemiology}
Heat stroke occurs most commonly during heatwaves, and heatwaves cause more deaths in many countries than any other natural disaster. It affects people at both extremes of age, outdoor workers, athletes, people with chronic disease, and those who cannot access cooling. Classic heat stroke affects mainly older adults with impaired temperature regulation or inadequate cooling, while exertional heat stroke affects young, fit people who exercise in extreme heat. Mortality from heat stroke is high, historically 10--50\%, though much lower with prompt treatment.

\section*{Aetiology and Pathophysiology}
\begin{itemize}
    \item \textbf{Classic heat stroke:} prolonged exposure to hot, humid conditions in people unable to cool themselves, often older adults without air-conditioning
    \item \textbf{Exertional heat stroke:} heavy physical activity in hot conditions overwhelms heat production and cooling, affecting athletes, soldiers, and workers
    \item \textbf{Pathophysiology:} when sweating and skin vasodilation cannot shed heat fast enough, core temperature rises; above 40 degrees Celsius, cellular proteins denature, the inflammatory response spirals, and multiple organs fail
    \item \textbf{Contributing factors:} dehydration, obesity, alcohol, some medications, and chronic illness
    \item \textbf{The distinction from fever:} heat stroke is environmental and causes failure of cooling, not the immune response of fever
\end{itemize}

\section*{Clinical Features}
\begin{itemize}
    \item Core body temperature of 40 degrees Celsius or higher
    \item Altered mental state: confusion, agitation, delirium, seizures, or loss of consciousness
    \item Hot, dry skin, although sweating may persist in exertional heat stroke
    \item Rapid, strong pulse and rapid breathing
    \item Headache, dizziness, and nausea
    \item Reduced or absent urine output as kidneys are affected
    \item Symptoms may progress rapidly from heat exhaustion
\end{itemize}

\section*{Differential Diagnosis}
\begin{itemize}
    \item Severe infection and sepsis cause high fever and confusion
    \item Stroke and other neurological emergencies cause altered consciousness
    \item Drug intoxication, including stimulants and anticholinergics, raises temperature and alters behaviour
    \item Thyroid storm and neuroleptic malignant syndrome cause hyperthermia
    \item Hypoglycaemia causes collapse and confusion
    \item Heat exposure history and temperature measurement are central to the diagnosis
\end{itemize}

\section*{When to Seek Medical Care}
Heat stroke is always a medical emergency. Anyone suspected of having heat stroke must have cooling started immediately while emergency services are called.

\textbf{Contact your local emergency services for:}
\begin{itemize}
    \item Confusion, agitation, seizures, or unconsciousness after heat exposure
    \item Body temperature of 40 degrees Celsius or higher
    \item Hot, dry skin with an inability to sweat
    \item Collapse or inability to stand
    \item Any person at high risk who is severely unwell in the heat
    \item Failure to improve within 30 minutes of cooling and fluids
\end{itemize}

\section*{Diagnosis and Investigations}
\begin{itemize}
    \item \textbf{Clinical diagnosis:} made on the combination of heat exposure, high core temperature, and altered mental state
    \item \textbf{Core temperature:} measured with a rectal or other accurate core probe; ear and skin thermometers are unreliable in the heat
    \item \textbf{Blood tests:} electrolytes, kidney and liver function, muscle enzymes, and clotting studies
    \item \textbf{Urine tests:} detect kidney injury and dehydration
    \item \textbf{Monitoring:} cardiac, respiratory, and neurological monitoring in hospital
    \item \textbf{Imaging:} CT of the brain if neurological signs are prominent or unclear
\end{itemize}

\section*{Management}
\begin{itemize}
    \item \textbf{Immediate cooling:} the priority is rapid lowering of core temperature; use cold water immersion, ice packs to the neck, armpits, and groin, fans with misting, and cool wet sheets
    \item \textbf{Emergency care:} ambulance transfer with cooling continued en route
    \item \textbf{Intravenous fluids:} correct dehydration and support circulation
    \item \textbf{Organ support:} intensive care for kidney, liver, heart, and respiratory support
    \item \textbf{Seizure management:} anticonvulsant medication where needed
    \item \textbf{Monitoring:} temperature, electrolytes, and organ function are tracked until stable
    \item \textbf{Recovery:} gradual resumption of activity after full recovery, with avoidance of heat
\end{itemize}

\section*{Complications}
\begin{itemize}
    \item Brain injury, including permanent neurological deficits and cognitive impairment
    \item Acute kidney injury and kidney failure
    \item Liver injury and failure
    \item Rhabdomyolysis, with muscle breakdown and dangerous electrolyte imbalances
    \item Disseminated intravascular coagulation and bleeding
    \item Heart failure, arrhythmias, and respiratory failure
    \item Death, particularly when cooling is delayed
\end{itemize}

\section*{Prognosis and Outcomes}
The outcome of heat stroke depends critically on how quickly body temperature is lowered. With immediate cooling and modern intensive care, most previously healthy people survive with full recovery. Delayed treatment, very high temperatures, and pre-existing disease worsen the outlook, and some survivors have lasting brain or kidney damage. Exertional heat stroke in young, fit people generally carries a better prognosis than classic heat stroke in the elderly. Prevention remains far more effective than treatment.

\section*{Prevention and Self-Care}
\begin{itemize}
    \item Never underestimate the heat: take heatwave warnings seriously
    \item Stay in air-conditioning or cooled spaces during extreme heat
    \item Drink water regularly and avoid alcohol and caffeine in the heat
    \item Schedule activity for early morning or evening and rest frequently
    \item Wear light, loose clothing, hats, and sunscreen
    \item Never leave children or older people in parked cars
    \item Check on vulnerable people regularly during heatwaves
    \item Learn and practise heat stroke first aid so you can act immediately
\end{itemize}

\section*{Sources and Further Reading}
The following reputable sources provide further information for healthcare students and patients seeking reliable, current advice about heat stroke.

\begin{itemize}
    \item Healthdirect. \textit{Heat stroke: first aid and treatment.} https://www.healthdirect.gov.au/heat-stroke
    \item St John Ambulance Australia. \textit{First aid for heat stroke.} https://stjohn.org.au/
    \item Better Health Channel. \textit{Heat stress and heat-related illness.} https://www.betterhealth.vic.gov.au/
    \item NHS. \textit{Heatstroke: symptoms and treatment.} https://www.nhs.uk/conditions/heat-exhaustion-heatstroke/
    \item Mayo Clinic. \textit{Heatstroke: first aid.} https://www.mayoclinic.org/
    \item World Health Organization. \textit{Heat and health.} https://www.who.int/
\end{itemize}
